\documentclass[11pt]{article}
\usepackage[T1]{fontenc}
\usepackage{lmodern}
\usepackage[letterpaper,margin=1in]{geometry}
\usepackage{amsmath,amssymb,amsthm,mathtools}
\usepackage{booktabs,array,longtable}
\usepackage{microtype}
\usepackage{xcolor}
\usepackage{listings}
\usepackage{enumitem}
\usepackage[hidelinks]{hyperref}
\usepackage{fancyhdr}
\hypersetup{pdftitle={Exact acceleration of the supplied unknot recognizer},
 pdfauthor={OpenAI},pdfsubject={Implementation, correctness arguments, and reproducible benchmarks}}
\pagestyle{fancy}\fancyhf{}
\fancyhead[L]{\small Exact acceleration of unknot recognition}
\fancyhead[R]{\small fastunknot 0.2.0}
\fancyfoot[C]{\thepage}
\setlength{\headheight}{14pt}
\setlength{\parskip}{3pt}
\setlist{nosep,leftmargin=1.6em}
\lstset{basicstyle=\ttfamily\small,breaklines=true,columns=fullflexible,
 keepspaces=true,showstringspaces=false,frame=single,framerule=0.3pt,
 aboveskip=8pt,belowskip=8pt}
\newtheorem{theorem}{Theorem}[section]
\newtheorem{proposition}[theorem]{Proposition}
\newtheorem{lemma}[theorem]{Lemma}
\newtheorem{corollary}[theorem]{Corollary}
\theoremstyle{definition}\newtheorem{definition}[theorem]{Definition}
\theoremstyle{remark}\newtheorem{remark}[theorem]{Remark}
\newcommand{\F}{\mathbb F_2}
\newcommand{\Kh}{\operatorname{Kh}}
\newcommand{\rKh}{\widetilde{\operatorname{Kh}}}
\newcommand{\rank}{\operatorname{rank}}
\newcommand{\poly}{\operatorname{poly}}
\newcommand{\code}[1]{\texttt{\detokenize{#1}}}
\newcommand{\UNK}{\textsf{UNKNOT}}
\newcommand{\KNOT}{\textsf{KNOTTED}}
\newcommand{\UNKNOWN}{\textsf{UNKNOWN}}

\title{Exact acceleration of the supplied unknot recognizer\\[5pt]
\large Sparse cancellation, bounded Jones obstructions,\\
and visible connected-sum factorization}
\author{Implementation and analysis report}
\date{18 September 2026 \quad---\quad Version 0.2.0}

\begin{document}
\maketitle
\begin{abstract}
We modify the actual \code{fast/} implementation in the submitted
\code{Knots.zip}, retaining its exact characteristic-two Khovanov backend.
The delivered software introduces fill-sensitive pivot scheduling, algebraic
memoization, a self-inverse-unit shortcut, a bounded scalar Jones state sum,
and diagrammatic connected-sum factorization. It also accelerates two
preprocessing routines and repairs crossing-sign handling for general PD
representatives. On same-host, fresh-process benchmarks, default recognition
of the Conway and Kinoshita--Terasaka fixtures is approximately 39 and 44 times
faster, respectively. A difficult 36-crossing raw scan finishes in 1.285 seconds,
whereas the unchanged baseline reaches a 20-second cap. An ablation using the
original uncached algebra confirms its reduced rank of 2949.

No unrestricted quasi-polynomial bound is established. We prove instead the
bound $\poly(n)+n2^{O(m)}$, where $m$ is the largest visible summand left by the
implemented factorization. This gives $n^{O(\log n)}$ on the class
$m=O((\log n)^2)$, and polynomial time on visible connected sums of
bounded-size diagrams. For a concrete family, it avoids the baseline's
exponentially large homology basis even when both polynomial-invariant
filters are disabled. The report separates these theorems from empirical
speedups, describes the exactness and resource-limit contracts, and includes
reproducible code, tests, inputs, and raw measurements.
\end{abstract}

\newpage
{\small\linespread{0.9}\selectfont\tableofcontents}
\newpage

\section{Scope, baseline, and the requested target}
\label{sec:scope}

The comparison target is the seven-module Python package under
\code{fast/fastunknot/} in the uploaded archive. It is preserved, byte for byte,
as \code{baseline/fastunknot/} in the deliverable. The modified package is
\code{fastunknot/}. No alternative implementation from the archive's six older
attempts is substituted as the baseline. Source provenance, including the
identity of the submitted ZIP, is recorded in \code{PROVENANCE.md}.

The baseline validates a classical, single-component planar diagram, performs
crossing-decreasing Reidemeister I and II moves, tests for a descending
basepoint, computes an Alexander polynomial, and finally computes Khovanov
rank over $\F$. Its last stage is a local scanning implementation in the style
of Bar-Natan \cite{BN}, not the naive full cube of resolutions. Our changes
build on that existing work; neither scanning nor the underlying knot
invariants are claimed as new discoveries.

\paragraph{What the supplied Lackenby materials say.}
The submitted February 2021 talk states the $n^{O(\log n)}$ target on physical
page 13 and outlines accelerated hierarchy operations on page 109
\cite{Ltalk}. The related July 2026 paper proves a hierarchy-based recognition
algorithm. Section 9 gives an iteration estimate in terms of hierarchy length
and surface complexity, rather than a crossing-number running-time bound
\cite{L2026}. Section~\ref{sec:target} explains precisely why our implementation
has not converted those operations into an unrestricted quasi-polynomial
algorithm.

There are distinct public slide versions: the author's March 2021 Oxford
variant states $2^{O((\log n)^3)}$, whereas the February file actually supplied
here states $n^{O(\log n)}$. We use the latter as the requested target rather
than silently identifying these different bounds \cite{Ltalk,Loxford}.

\paragraph{The delivered result.}
The software is executable using Python 3.10 or later and its standard library.
It has no missing hierarchy subroutine hidden behind an oracle. With unlimited
resources it retains an exact complete recognition route; configured limits
can instead return \UNKNOWN. Its general worst-case bound remains exponential.
The strongest new asymptotic statement is a visible-factor parameter bound,
and the strongest observed backend improvement comes from changing the
cancellation order, not merely adding a faster knot invariant.

\section{Representation and trusted mathematical ingredients}
\label{sec:representation}

A PD row $(a,b,c,d)$ lists the four ports of a crossing counterclockwise.
Ports 0 and 2 belong to the underpassing strand, and ports 1 and 3 to the
overpassing strand. Each edge label occurs twice. The public constructors
normalize labels, check the opposite-port traversal has one component, and
check the rotation system is spherical. The empty diagram denotes one unknot,
not an empty link. These conventions are inherited and matter for every
subsequent exactness statement.

\subsection{The scanning category}

At a partial stage, a crossing-free object is a matching $M$ of the current
boundary endpoints, together with a homological degree $h$. Closed circles
are immediately delooped into two copies, labeled $1$ and $x$. A morphism
between matchings is a characteristic-two sum of dotted surfaces, represented
in a basis indexed by dot subsets of the circles formed by joining the two
matchings. The relevant Frobenius algebra is
\[
 A=\F[x]/(x^2),\qquad
 \Delta(1)=1\otimes x+x\otimes1,\quad
 \Delta(x)=x\otimes x,\quad
 \epsilon(1)=0,\quad\epsilon(x)=1.
\]
The baseline already implements tensoring with a crossing, delooping,
composition, and cancellation of invertible differential entries. These are
local versions of the formal-complex operations of \cite{BN}. Our optimized
scanner retains the same mathematical computation.

At the end of a closed diagram, the matching is empty. All nonzero
characteristic-two scalar differential entries are units, so cancellation
leaves a zero differential. The remaining number of objects is the total
unreduced Khovanov rank. Homological degrees retained by this implementation
are \emph{raw cube degrees}; writhe shifts and quantum gradings are not
reported. Total ranks suffice for recognition.

\subsection{Why unreduced rank two detects the unknot}

We use the established theorem that reduced Khovanov homology of rank one
detects the unknot \cite{KM}. Its application to the particular field and
unreduced implementation deserves explanation.

\begin{lemma}\label{lem:double}
For a knot diagram, after forgetting the quantum grading,
\[
 \dim_{\F}\Kh^h(K;\F)=2\dim_{\F}\rKh^h(K;\F)
\]
in each homological degree.
\end{lemma}
\begin{proof}
Let $X$ multiply the label on a marked circle by $x$. Let $\nu$ be the sum of
the operations that change one circle label $x$ to $1$ and annihilate label
$1$. The formulas for $m$ and $\Delta$ show that $\nu$ commutes with the
characteristic-two differential. They also give
\[
 X^2=\nu^2=0,\qquad X\nu+\nu X=1.
\]
Consequently $X\nu$ and $\nu X$ are complementary chain projections.
Their images are isomorphic by $X$ and $\nu$, and the first is the reduced
subcomplex $\operatorname{im}X$. Neither operator changes homological degree.
Taking homology gives the assertion.
\end{proof}

If the reduced $\F$ rank is one, the universal coefficient theorem implies
that the reduced rational rank is at most one. It is at least one: the
normalized reduced Euler characteristic has evaluation one at the unit
parameter, so the rational complex cannot have zero homology. The
rank-one unknot-detection theorem therefore applies. Conversely the reduced
homology of the unknot has rank one. Thus the backend's unreduced rank test
is exactly a rank-two test, not a heuristic comparison with a familiar knot
table. The software does not reprove the deep detection theorem; it relies on
it just as the baseline does.

\section{Accelerating the categorical scanner}
\label{sec:scanner}

\subsection{Fill-sensitive cancellation}

Write a differential block around an invertible entry $\phi:B\to C$ as
\[
 d=\begin{pmatrix}\phi&\delta\\ \gamma&\varepsilon\end{pmatrix}
 :B\oplus U\longrightarrow C\oplus V.
\]
Removing the contractible pair $B,C$ replaces the remaining block by
\begin{equation}\label{eq:schur}
 \varepsilon'=\varepsilon+\gamma\phi^{-1}\delta.
\end{equation}
The plus sign is correct over $\F$. This cancellation is a chain-homotopy
equivalence; it is the categorical Gaussian elimination used in the baseline
and in \cite{BN}. The algebraic validity does not depend on which available
invertible entry is chosen.

For a candidate arrow $b\to c$, let $r_c$ be the number of arrows entering $c$
and $s_b$ the number leaving $b$. Cancelling it may update at most
\begin{equation}\label{eq:score}
  (r_c-1)(s_b-1)
\end{equation}
other entries. The old implementation uses a stack of candidates. The new
one uses a minimum heap keyed by this local fill estimate, then by a serial
number for deterministic ties. A candidate is rechecked for existence and
invertibility before cancellation. Its numerical score is the score at
insertion, not a continuously maintained exact minimum.

\begin{proposition}
Stale scores in this queue cannot invalidate the homology computation.
\end{proposition}
\begin{proof}
A score only selects which candidate to inspect. The current morphism and its
current source and target matchings are checked again. Every performed
operation is therefore a valid instance of~\eqref{eq:schur}. Removed or no
longer invertible candidates are discarded. Any entry that changes to a
nonzero unit is reinserted. The process terminates because each cancellation
removes two objects. Selection by a stale score changes performance, not the
resulting homotopy type.
\end{proof}

This is a heuristic for sparsity, not a theorem that the number of intermediate
objects is polynomial. It is nevertheless the decisive change on the
36-crossing fixture. A separate ablation changes only this cancellation loop
while continuing to use the baseline's uncached algebra.

\subsection{Units are self-inverse in the implemented algebra}

For a matching with $k$ arcs, its endomorphism ring in this normal form is
\begin{equation}\label{eq:ring}
 R_k=\F[x_1,\ldots,x_k]/(x_1^2,\ldots,x_k^2).
\end{equation}
The basis monomials are square-free dot subsets. The identity is the empty
subset, and an endomorphism is a unit exactly when the identity coefficient
is one.

\begin{lemma}\label{lem:inverse}
Every unit of $R_k$ is its own inverse.
\end{lemma}
\begin{proof}
Write a unit as $1+u$, where $u$ has zero constant coefficient. In a commutative
ring of characteristic two the square of a sum is the sum of the squares.
Every nonconstant square-free monomial has square zero in~\eqref{eq:ring}.
Thus $u^2=0$ and $(1+u)^2=1$.
\end{proof}

The old inverse routine computes a nilpotent expansion using general surface
composition. The replacement returns a copy of the input morphism. The
restriction to \emph{endomorphisms with identical source and target matching}
is essential: a term that looks empty in a general Hom space is not by itself
an identity between different objects. The pivot and identity fast paths both
check this restriction. Exhaustive tests enumerate all 128 units for a
three-arc matching and square each using the original composition routine.

\subsection{Memoization without mutable-alias errors}

The same small matching and the same local morphism can occur on many copies
of a chain object. Recomputing their topological gluing for each copy wastes
work unrelated to the number of genuinely different local maps.

The modified implementation has three kinds of reuse. First, composition
returns zero immediately on a zero input and copies the other argument for a
valid identity input. Otherwise its key contains the three matchings and
\emph{frozen} copies of the two morphisms. The cached value is frozen, and the
public result is copied back into an owned mutable set. Second, geometric
circle and smoothing-gluing caches are bounded to 8192 entries each; the
composition cache is bounded to 32768. Third, each crossing has a temporary
cache of the delooped entries for a given pair of matchings, morphism, and
pair of smoothing indices. Boundary and crossing data are constant within
that temporary cache, which is discarded at the end of the crossing.

The entry maps shared between chain arrows are read-only by convention.
When cancellation updates a differential entry it first copies it and then
performs symmetric difference. This prevents one arrow's cancellation from
mutating another arrow or a cached result. The cache limits count entries,
not bytes; a single morphism can still be large. \code{clear_caches()} permits
explicit release between applications.

On the recorded difficult scan, the crossing-map cache has 195674 hits and
3154 misses. The optimized run performs 16884 cancellations and records
168172 calls to composition inside cancellation, with a maximum of 17694
objects before elimination. These are measured work counters, not asymptotic
bounds.

\subsection{Exactly the same greedy order, faster}

For an unprocessed crossing, write $s$ for the number of its ports incident
to the current boundary, $b$ for the boundary size, and $\ell$ for its number
of repeated edge labels. The baseline's lexicographic score is
\[
 (s,-b-4+2s+2\ell,-i),
\]
where $i$ is the crossing index. Since $b$ is common to all candidates and $s$
is compared first, this induces exactly the same order as $(s,\ell,-i)$.
Only neighbors of a newly scanned crossing can change $s$. On a four-valent
PD graph there are $O(n)$ such updates in total.

Maintaining this score in a versioned heap therefore requires
$O(n\log n)$ heap time and $O(n)$ heap records per starting crossing, instead
of rescoring $O(n)$ candidates at $O(n)$ stages. The multi-start routine
retains the baseline's start selection and boundary-profile tie breaking.
Its default samples only a bounded number of starts, hence remains
$O(n\log n)$ overall. Explicitly requesting every start instead costs
$O(n^2\log n)$. Tests compare the actual returned permutations with the
baseline for every start on 100 generated diagrams.

\section{A bounded, exact Jones obstruction}
\label{sec:jones}

The Jones layer is a separate scalar computation. It does not call the
categorical surface-composition or gluing routines and does not compute
Khovanov homology. The shared infrastructure consists only of validated
diagrams, orientation, and a crossing-order heuristic.

\subsection{State sum and finite-ring specialization}

Use the bracket convention $\langle\bigcirc\rangle=1$, smoothings
$(a,b)(c,d)$ and $(a,d)(b,c)$ with respective weights $A$ and $A^{-1}$, and
\[
 \delta=-A^2-A^{-2}.
\]
For a nonempty diagram $D$ with $n$ crossings, a smoothing state $s$ with
$|s|$ second smoothings and $c(s)$ circles contributes
\begin{equation}\label{eq:bracket}
 \langle D\rangle=
 \sum_{s\in\{0,1\}^n}A^{n-2|s|}\delta^{c(s)-1}.
\end{equation}
The writhe-normalized quantity is
\begin{equation}\label{eq:jones}
 J_A(D)=(-A^3)^{-w(D)}\langle D\rangle.
\end{equation}
It is the usual normalized bracket version of the Jones invariant, with
variable choice $t=A^{-4}$ \cite{Kauffman}. Our PD and writhe conventions
are paired consistently; swapping to a mirror replaces the corresponding
variable convention but does not change the unknot criterion.

The implementation evaluates in $\mathbb Z/M\mathbb Z$. It checks that both
$A$ and $\delta$ are invertible and uses exact modular inverses. Primality of
$M$ is not needed. The default is $M=1000000007$, $A=2$.

\begin{proposition}\label{prop:jones}
A successfully computed residue $J_A(D)\ne1$ certifies \KNOT. The residue
$J_A(D)=1$ gives no acceptance criterion for \UNK.
\end{proposition}
\begin{proof}
For an unknot, the normalized Laurent polynomial~\eqref{eq:jones} is
identically one. Specializing it at units in any admissible finite ring still
gives one. Contraposition proves the rejection criterion. Distinct Laurent
polynomials may agree at the chosen ring element, so equality at one value
cannot prove polynomial equality or unknotness.
\end{proof}

This is not a randomized test with a small probability of false rejection.
A collision can only make the filter unhelpful. No conjecture about whether
the Jones polynomial detects the unknot is used.

\subsection{Frontier dynamic programming}

After processing some crossings, store a dictionary from boundary matchings
to scalar partial weights. Adding a crossing has two transitions per matching.
A small disjoint-set computation joins the old matching to the chosen local
smoothing. A resulting component has either two new boundary endpoints,
which become an arc in the new matching, or no endpoints, which contributes
one factor $\delta$. Anything else is rejected as inconsistent geometry.
Weights of equal new matchings are summed modulo $M$.

For convenience, the dynamic program accumulates $\delta^{c(s)}$ rather than
$\delta^{c(s)-1}$; it divides by $\delta$ once at the end. The zero-crossing
unknot is treated separately as value one. Induction on processed crossings
proves that the dictionary is exactly the partial sum obtained by aggregating
all smoothing choices with the same boundary matching. Final closure then
gives~\eqref{eq:bracket}, and writhe normalization gives~\eqref{eq:jones}.

If the boundary has at most $w$ endpoints, the number of arbitrary matching
states is at most $(w-1)!!=2^{O(w\log w)}$. We do \emph{not} substitute a
Catalan bound without proving that the boundary lies on one disk in the
appropriate cyclic order: a general scan need not have that form. There are
also at most $2^n$ partial smoothing states before aggregation. This scalar
algorithm can be exponential and is not the complete recognizer.

The default filter is bounded by 100000 stored states and 200000 transitions.
Reaching either work cap skips the filter and continues with Alexander and
then the exact scanner. With a fixed-size modulus and the default work cap,
its overhead is polynomial in input size. A global time deadline has a
different meaning: it terminates the recognition attempt with \UNKNOWN.

\subsection{Orientation repair needed for exact normalization}

The baseline's crossing signs depend only on which over-port is incoming,
as though the incoming under-port were always slot 0. A legal two-slot
rotation of an individual PD row preserves the crossing but can violate that
assumption. Let $u\in\{0,2\}$ and $o\in\{1,3\}$ be the actual incoming
under- and over-port. The repaired local sign is
\[
 \varepsilon=\begin{cases}+1,&o-u\equiv3\pmod4,\\-1,&o-u\equiv1\pmod4.
 \end{cases}
\]
Simultaneously adding two to both incoming ports leaves it unchanged.

It would be unsafe to change \code{signs()} alone: Alexander matrix
construction previously used fixed under-port positions in conjunction with
the old effective sign. It now identifies the actual incoming under-arc and
interchanges its under-arc variables when needed. Tests check writhe,
Alexander polynomial, and Jones residues under independent two-slot row
rotations on 100 generated diagrams, as well as reversal of writhe on mirrors.

\section{Visible connected sums: avoiding a tensor-product explosion}
\label{sec:factors}

The baseline always computes a final homology basis if its Alexander test
fails. A small frontier cannot prevent this basis from being enormous.
Connected sums expose a way to avoid constructing it while still deciding
unknotness, or even computing its rank.

\subsection{Finding and replaying a cut}

Traverse the oriented knot once and record the crossing index at each visit.
The resulting cyclic Gauss word has length $2n$, with each crossing appearing
twice. For each possible start, extend a proper cyclic interval and maintain
the set of crossing indices seen an odd number of times. An empty odd set
means that every crossing touched by that interval has both visits inside it.

\begin{lemma}\label{lem:cut}
A nonempty proper closed interval of this kind determines a visible
connected-sum decomposition of the spherical knot diagram.
\end{lemma}
\begin{proof}
Let $S$ be the set of its crossings. Along the knot traversal, every edge
between visits inside the interval has both endpoints in $S$, and similarly
for the complementary interval. Precisely the incoming and outgoing boundary
edges cross from $S$ to its complement. Both induced subgraphs are connected,
since each is traversed by its interval. Thus the two boundary edges form a
bond of the plane projection graph. Its dual is a simple two-edge cycle,
realizable as a Jordan curve avoiding vertices and meeting those two edges
once each. This curve is a diagrammatic connected-sum circle. Closing each
of the two resulting arcs gives the two summands.
\end{proof}

The implementation closes a side by identifying its two cut labels, retains
its PD rows, and runs the usual diagram validator. It repeats this operation
on an explicit stack, avoiding recursion-depth limits. A search on a
$k$-crossing diagram costs $O(k^2)$ set operations; at most $n-1$ cuts give
the conservative overall bound $O(n^3)$, before label bit costs. A more
specialized linear-time bond algorithm could reduce this overhead, but is
not needed for the claimed bound.

Each cut records its current node number, child numbers, the two edge labels,
and the chosen crossing subset in that node's normalized PD. The replay
routine checks the exact cut edges, reconstructs both closures, and validates
the resulting diagrams. These records describe actual visible cuts, not
asserted knot names. Failure to find a cut says nothing about topological
primality: a composite knot can have a projection without such an exposed cut.

\subsection{Reduced ranks multiply}

\begin{proposition}\label{prop:tensor}
For the visible connected sum $K\#L$ over $\F$, reduced Khovanov complexes
have the connected-sum tensor description, and hence
\begin{align}
 r(K\#L)&=r(K)r(L),\label{eq:rankproduct}\\
 r_h(K\#L)&=\sum_{i+j=h}r_i(K)r_j(L),\label{eq:degreeproduct}
\end{align}
where $r$ is total reduced rank and the $r_h$ use the compatible raw cube
homological degrees of the spliced diagrams.
\end{proposition}
\begin{proof}
Mark a point on each spliced edge. Each full cube state is a free $A$-module
through the marked circle. A state of the connected sum consists of one state
of each summand, with the marked circles identified; its module is their
tensor product over $A$. The crossing differentials lie in one summand or
the other, and this identification respects multiplication and
comultiplication. Thus the connected-sum complex is the corresponding tensor
product over $A$.

Reduced complexes may be represented by base change along $A\to A/(x)=\F$;
this is isomorphic, up to the irrelevant quantum shift, to taking marked
label $x$. Base change of the displayed tensor product yields the tensor
product of the two reduced complexes over $\F$. Homology of a tensor product
of finite complexes over a field is the tensor product of their homologies.
Its dimension and homological grading give~\eqref{eq:rankproduct}
and~\eqref{eq:degreeproduct}.
\end{proof}

Every summand has positive integral reduced rank. Therefore the product is
one exactly when every factor is one. Together with the detection theorem,
this gives a particularly simple exact decision strategy: accept only when
all visible factors are accepted, and reject as soon as one is proved knotted.
It does not require a general topological prime-factorization algorithm.

The separate \code{factorized_khovanov_rank} function computes the product
and degree convolution instead of a basis. Lemma~\ref{lem:double} converts
the scanner's per-degree unreduced counts into reduced counts by dividing
by two, then doubles the final counts for the output's unreduced convention.
An odd intermediate count is treated as an error. Identical normalized PD
leaves are memoized; this is an exact equality test, not a knot-isomorphism
heuristic. Different presentations of the same knot may be recomputed.

\subsection{An exponential separation from the supplied algorithm}
\label{sec:family}

Let $D$ be the supplied eleven-crossing Conway fixture and form $D^{\# k}$
by the edge splice in \code{tools/common.py}. On $D$, the independent dense
reference computation gives reduced rank $33$, the exact Alexander routine
gives $1$, and inspection finds no legal R1 or R2 move. These checks are
retained in \code{results/independent-dense-named.json}.

\begin{proposition}\label{prop:separation}
On this family, the supplied complete baseline requires
$\Omega(33^k)$ stored final chain objects. The new recognition algorithm
runs in polynomial time in $k$, even with both Jones and Alexander filters
disabled.
\end{proposition}
\begin{proof}
The repeated edge splice displays all $k$ factors. It introduces no crossing,
no monogon, and no new two-crossing bigon: the two splicing edges have
distinct endpoints in each copy, and a newly joined face runs through both
sides of the splice. Old faces not incident to the splice are unchanged.
Consequently the baseline's crossing-decreasing reductions do not remove
these factors. The connected sum is nontrivial by~\eqref{eq:rankproduct},
so it has no descending basepoint. Its Alexander polynomial is the product
of $k$ unit polynomials, still one.

Thus the baseline reaches its scanner. The final unreduced homology has
rank $2\cdot33^k$. That implementation has zero final differential and
explicitly retains one object per generator, hence must store at least that
many objects before returning. This lower bound is on this supplied
implementation, not on every possible algorithm for the decision problem.

The new decomposition finds visible factors of size eleven in polynomial
time. It computes the first such factor's rank and immediately rejects;
there is no need to construct any tensor-product basis. The fixed-size
factor computation is a constant in the family parameter $k$. Polynomial
decomposition and preprocessing therefore dominate.
\end{proof}

The same argument applies to any fixed reduced diagram of a knot with
Alexander polynomial one and reduced rank greater than one, with an exposed
splice that preserves the absence of reducing moves. It is a decision-specific
improvement: building the whole homology is unnecessary for deciding whether
its reduced dimension equals one.

The factored-rank API goes further on the concrete $k=12$ fixture: it returns
\[
 r(D^{\#12})=33^{12}=1667889514952984961
\]
in about 0.032 seconds in the recorded run. The integer has a small binary
encoding, although a vector-space basis with that many elements does not.
The homological count array is also compact. This measurement is for the
factored-rank algorithm, not for a raw scan of 132 crossings.

\section{The complete recognition pipeline}
\label{sec:pipeline}

The new default is
\[
\begin{gathered}
\text{validate}\ \longrightarrow\ \text{R1/R2 reduction}
\ \longrightarrow\ \text{descending test}\\
\longrightarrow\ \text{visible factorization}
\ \longrightarrow\ \text{per-factor bounded Jones test}\\
\longrightarrow\ \text{per-factor Alexander test}
\ \longrightarrow\ \text{accelerated exact scan}.
\end{gathered}
\]
Each factor is simplified and tested for descent again. Small factors are
processed first. Repeated identical PD leaves share their computed decision.
A knotted factor ends the whole computation immediately; all factors must
be trivial for acceptance. A factor whose scanner hits an object ceiling is
unknown, but another factor can still prove knottedness unless the global
time budget itself has expired.

\subsection{The linear descending test}

For a cyclic traversal, let $B$ be the number of crossings first encountered
underneath from a chosen basepoint. The diagram is descending exactly when
$B=0$. Moving the basepoint past one visit reverses the first-encounter status
of only that crossing. Its contribution changes by $+1$ if the passed visit
was over, and by $-1$ if it was under. Compute $B$ at one start in linear time,
then sweep all starts with constant work per visit. Repeating for the opposite
orientation gives an $O(n)$ procedure. The baseline recomputes first visits
for each possible start, taking $O(n^2)$ time.

A descending diagram is an unknot: after cutting at the basepoint, height can
be assigned monotonically in traversal order, and the resulting long arc can
be straightened before reconnecting its ends. Failure to find a descending
start is not a negative certificate. Our implementation returns the actual
successful dart and preserves this one-sided use.

\begin{theorem}\label{thm:soundness}
On a validated classical knot diagram, every completed decision of the new
recognizer is mathematically sound, assuming the correctness of the stated
invariant theorems and the implemented exact local operations. With no
resource ceilings and sufficient memory it terminates with a decision.
\end{theorem}
\begin{proof}
R1 and R2 moves preserve knot type and strictly decrease crossing number.
The descending test only accepts trivial knots. The cut routine produces
connected-sum decompositions by Lemma~\ref{lem:cut}. The Jones residue only
rejects by Proposition~\ref{prop:jones}; Alexander rejects only a polynomial
different from its normalized unknot value. Neither filter's failure accepts
an unknot.

An unresolved factor is sent to the same finite crossing-complex construction
as in the baseline. Every modified cancellation is a chain-homotopy
equivalence, and the caches only reuse exactly equal local computations.
Its total rank gives the exact decision described in
Section~\ref{sec:representation}. Equations~\eqref{eq:rankproduct} and
\eqref{eq:degreeproduct} justify combining summand decisions. All loops are
finite on an unlimited run, so an unresolved factor cannot be left behind.
Resource exceptions instead yield \UNKNOWN and are not used as evidence for
either mathematical verdict.
\end{proof}

The qualification concerning implemented local operations is not a claim of
formal verification. Ordinary source code can have defects; the independent
comparisons and algebra tests in Section~\ref{sec:validation} provide
substantial empirical checks, not a machine-checked proof of the Python
interpreter or of the whole implementation.

\section{Complexity: what is proved and what is not}
\label{sec:complexity}

\subsection{An unrestricted exponential upper bound}

Let $n$ be the number of crossings in a normalized PD diagram. At one crossing,
an old matching has two smoothing choices and each creates at most two new
closed circles. Delooping therefore makes at most $2\cdot2^2=8$ objects from
one old object. Before reductions, the number of objects in any stage is
bounded by $8^n$. Cancellation never increases it.

A Hom-space normal form has $O(n)$ possible circle keys, so a morphism has
at most $2^{O(n)}$ square-free terms. Geometry construction, composition,
differential storage, and an individual Schur update are polynomial in the
explicit object count and morphism-support sizes. There are at most a
polynomial number of candidate updates in the explicit stage size: each
cancellation removes objects, and each one touches at most the square of the
remaining number of objects. The heap can have duplicate candidates, but at
most a polynomial number in these same sizes. Consequently these deliberately
loose estimates still give
\begin{equation}\label{eq:exp}
 T_{\rm scan}(n)=2^{O(n)}
\end{equation}
including polynomial-size label and object-index arithmetic. This is an
upper bound, not an empirical model predicting a useful base for the
exponential.

The Alexander routine is fraction-free elimination over integer polynomials.
Degrees of relevant minors and bit sizes of their coefficients are polynomial
in diagram size; polynomial arithmetic and pivoting therefore add polynomial
time. Initial canonicalization adds a polynomial in the actual input bit
length when external labels have long encodings. The ordinary statement in
$n$ assumes normalized labels of $O(\log n)$ bits, as supplied to the backend
by the validated constructors.

\subsection{The visible-factor parameter}

\begin{definition}
Let $m(D)$ be the maximum crossing count of a leaf returned by the implemented
visible decomposition after initial simplification. For the already accepted
empty diagram, set $m(D)=0$. This is an algorithmically produced parameter of
the input projection, not the crossing number of the largest topological
prime factor.
\end{definition}

\begin{theorem}\label{thm:parameter}
With default bounded Jones filtering and no terminating resource limit, the
recognizer has the bound
\begin{equation}\label{eq:parameter}
 T(D)=\poly(n)+n2^{O(m(D))}.
\end{equation}
The optional factored-rank computation satisfies the same form of bound with
$m$ instead defined as its largest \emph{raw-input} decomposition leaf, since
that API does not initially simplify. Integer rank arithmetic and degree
convolution add polynomial overhead.
\end{theorem}
\begin{proof}
Simplification, the descending sweeps, and repeated cut searches take
polynomial time. There are at most $n$ nonempty leaves. Each has at most
$m=m(D)$ crossings, and further reduction cannot increase this number.
The bounded Jones filter contributes polynomial overhead; exact Alexander
computation is polynomial in its factor size. Bound~\eqref{eq:exp} on every
remaining factor gives at most $n2^{O(m)}$ scanning work. Early rejection or
memoization can only reduce this bound.

For complete factorized ranks, apply the scanning estimate to the raw
decomposition leaves of that API. A raw degree lies between zero and the total
crossing count. A count is at most the full cube dimension, hence has $O(n)$
bits by the same exponential dimension bound. Products and successive
convolutions of the $O(n)$-length degree arrays therefore cost polynomial
time in $n$. The large numerical value of a rank does not require an equally
large number of storage cells when it is returned as an integer.
\end{proof}

\begin{corollary}\label{cor:qp}
The delivered recognizer is polynomial-time on inputs with $m=O(\log n)$,
and has the requested $n^{O(\log n)}$ bound on inputs with
$m=O((\log n)^2)$. In particular, visible connected sums of bounded-size
diagrams are polynomial-time inputs.
\end{corollary}
\begin{proof}
Substitute the respective bounds into~\eqref{eq:parameter}, using
$2^{O((\log n)^2)}=n^{O(\log n)}$.
\end{proof}

The condition on $m$ is an additional hypothesis. There is no argument that
all knot diagrams have small $m$, and no such claim is made in the code or
its JSON output. Without that condition, $m$ can equal $n$.

\subsection{Why a thin boundary does not prove fast Khovanov rank}

The baseline README suggests that its exponential part depends only on
boundary size. That statement overlooks multiplicities. A constant-width
scan of repeated bounded-size summands still ends with $2\cdot33^k$
objects for the family in Section~\ref{sec:family}. There are only finitely
many matching types at bounded boundary size, but arbitrarily many copies
of those types in a chain complex. The final boundary is even empty.

A width-only count of matching types therefore cannot bound this
implementation's full homology computation. This is not an objection to
local scanning as a practical algorithm. It identifies the exact quantity
that must be compressed or avoided: multiplicities and their differentials,
not just topological types of frontier states. Visible factorization solves
that problem for an exposed tensor-product structure, not for arbitrary
prime diagrams.

\subsection{Attempting to reach the hierarchy target}
\label{sec:target}

Lackenby's Section 9 analyzes a hierarchy with maximal length $L$ and maximal
surface pattern-complexity $g$. Encoding the complexity sequence as base
$g+1$ digits bounds the number of step executions by
\begin{equation}\label{eq:hierarchy}
 L(g+1)^L.
\end{equation}
That section explicitly does not provide the required bounds on both
parameters or precise costs for all the steps \cite{L2026}. It cannot simply
be read as a proved running time in the crossing number.

Here is the conditional calculation relevant to the requested target. If an
implemented hierarchy representation has polynomially bounded $g$, depth
$L=O(\log n)$, polynomial-size intermediate encodings, and polynomial-cost
operations, then~\eqref{eq:hierarchy} is $n^{O(\log n)}$. Removing any of these
premises leaves a gap. A short list of abstract operations or a lexicographic
termination argument alone does not establish the premises.

We examined whether the present scanner could supply this compression by
small frontier states. The multiplicity obstruction above defeats that route
in general. We instead implemented the concrete special case where the
projection exposes a tensor decomposition, obtaining
Theorem~\ref{thm:parameter}. The hierarchy-specific operations in the talk's
final diagram---hierarchical multisurface construction, Cheeger-region use,
Haken-lemma transformations, weak reduction, and multisurface
simplification---are not implemented here. No executable routine is
misrepresented as implementing them, and the optional limits are not
substitutes for their cost analysis.

This describes the boundary of the present result. It is neither a
refutation of the announced hierarchy approach nor a claim that such an
implementation cannot be developed. The delivered improvement is the
verified-in-experiments exact scanner and filter system, with the restricted
complexity theorem just proved.

\section{Reproducible experiments}
\label{sec:experiments}

\subsection{Protocol and machine}

The included \code{tools/benchmark.py} creates an isolated Python process for
each sample. Both backends are imported before the timer; JSON parsing and
construction and validation of the same diagram are also outside the timer.
All algorithmic work for the selected operation, including order selection
when used, is inside. There is no cache warming between samples. Completed
rows have three samples, summarized by their median. A capped row stops after
its first unsuccessful sample rather than pretending to have a completion
median.

Measurements were collected with Python 3.13.5, Linux x86-64 with glibc 2.41,
on a virtualized host identifying its processor as AMD EPYC 9V74. The supplied
baseline runs in that same Python and environment. The cooperative budget is
20 seconds; the parent process also has a 24-second emergency timeout.
All recorded capped rows in this run reached the cooperative limit. Neither
method is forced to a smaller object ceiling. The routine full benchmark does
not enable the expensive $d^2=0$ debug checks; correctness runs enable them
separately.

There are 37 backend/case/operation rows and 105 individual samples. Raw
JSON includes the full answers, orders, work counters, timing, and process
peak resident memory where available. These numbers are microbenchmarks of
the included fixtures, not estimates of performance on a random distribution
of knots. The benchmark is single-threaded and no external algebra package
or compiled accelerator is used.

\subsection{Default end-to-end algorithmic timings}

\begin{table}[htbp]
\centering
\begin{tabular}{@{}lrrr@{}}
\toprule
Fixture & Baseline (ms) & New (ms) & Ratio\\
\midrule
Conway, 11 crossings & 35.630 & 0.915 & 38.9\\
Kinoshita--Terasaka, 11 & 39.297 & 0.886 & 44.3\\
Hard unknot, 8 & 1.621 & 1.819 & 0.89\\
Five-strand example, 36 & 8.082 & 5.217 & 1.55\\
Unknot braid, 40 & 6.411 & 6.353 & 1.01\\
Torus knot $T(3,5)$, 10 & 0.590 & 0.904 & 0.65\\
\bottomrule
\end{tabular}
\caption{Default recognition. The ratio is baseline time divided by new time;
a value below one is a regression. Inputs and imports are excluded for both.}
\label{tab:pipeline}
\end{table}

Conway and Kinoshita--Terasaka both have Alexander polynomial one in these
fixtures, so the old pipeline reaches Khovanov homology. The new modular
Jones filter rejects them immediately. This is the source of the roughly
39-fold and 44-fold end-to-end improvements, not just a constant-factor
scanner optimization.

The new pipeline is \emph{not} universally faster. Extra exact filters add
overhead to the small hard unknot, which still needs positive recognition,
and to $T(3,5)$, which the old Alexander test already rejects cheaply. The
latter becomes approximately 1.53 times slower, although still submillisecond
in these measurements. Users with such a workload can disable Jones and/or
visible factorization. No general dominance claim is justified by the data.

\subsection{Raw scanning, without preprocessing or invariant shortcuts}

\begin{table}[htbp]
\centering
\begin{tabular}{@{}lrrr@{}}
\toprule
Fixture & Baseline (s) & New (s) & Reduced rank\\
\midrule
Conway & 0.033667 & 0.010819 & 33\\
Kinoshita--Terasaka & 0.039130 & 0.011889 & 33\\
Hard unknot, 8 & 0.002139 & 0.001898 & 1\\
Five-strand example, 36 & $>20$ (capped) & 1.284672 & 2949\\
Unknot braid, 40 & 0.008588 & 0.004003 & 1\\
Torus knot $T(3,5)$ & 0.010973 & 0.010467 & 7\\
\bottomrule
\end{tabular}
\caption{The same original PD is sent straight to each scanner. The capped
entry is not a measured runtime; the difficult-case speed ratio is greater
than 15.5 relative to the 20-second budget.}
\label{tab:scan}
\end{table}

The hard 36-crossing case already appears in the supplied benchmark data.
That older file records a 600-second raw-scan timeout under Windows and
Python 3.14.4. We do \emph{not} divide that time by a new Linux measurement.
The comparison in Table~\ref{tab:scan} reruns the actual original code on the
same host, with a clearly reported 20-second cap.

For this diagram, the reduced rank 2949 is new computational output relative
to that capped baseline row. The raw result has unreduced rank 5898 and
maximum boundary size eight. These figures illustrate both the practical
value and the limitation of narrow scans: a thin scan can still have
substantial homology multiplicities.

The default old pipeline already decides this example in approximately
8 milliseconds after six R2 moves reduce it to 24 crossings and Alexander
rejects it. The new pipeline uses a Jones residue of 426151971 modulo
1000000007, obtained with 14 maximum states and 400 transitions. Its total
time is approximately 5 milliseconds. Thus the large raw-scanner improvement
must not be presented as a comparable end-to-end improvement on this input.

\subsection{Ablating pivot policy and algebra caching}

\begin{table}[htbp]
\centering
\begin{tabular}{@{}p{3.45in}rr@{}}
\toprule
36-crossing raw-scan variant & Seconds & Reduced rank\\
\midrule
Original algebra, original stack policy & $>20$ & not completed\\
New caching, original stack policy & $>20$ & not completed\\
Original uncached algebra, fill policy only & 7.627084 & 2949\\
New caching and fill policy & 1.284672 & 2949\\
\bottomrule
\end{tabular}
\caption{Ablation results. The fill-only worker installs the new elimination
loop in the baseline module namespace, retaining the original composition,
inversion, crossing maps, and scan-order implementation.}
\label{tab:ablation}
\end{table}

The comparison shows that caching alone does not resolve this particular
bottleneck: its stack-policy variant also reaches the cap. Changing pivot
selection makes the original uncached algebra finish. Adding the cache and
unit shortcuts then improves that fill-only time by approximately 5.94.
The four rows do not isolate every small optimization individually, so this
is not a claim to have attributed a precise fraction of the speedup to each
cache.

The fill-only variant and the optimized variant agree on every reported
homological degree count, not just the final total. A separate fill-only run
with $d^2=0$ checks also completes and gives rank 2949; its result is saved in
\code{results/hard36-original-algebra-check.json}. This tests the optimized
algebra against the original one, but both variants still share the same
underlying categorical model. It is not an independent formal proof of that
model.

For context, the first successful optimized sample has process peak RSS
129708 KiB, whereas the first fill-only sample has 168744 KiB. These include
interpreter and process memory, not just live chain-complex storage. They
are not a proven memory bound and should not be extrapolated to arbitrary
knots or scan widths.

\subsection{Factorization and ordering experiments}

\begin{table}[htbp]
\centering
\begin{tabular}{@{}lrrr@{}}
\toprule
Operation & Baseline (s) & New (s) & Ratio\\
\midrule
Two Conway summands, forced scan & 1.427456 & 0.010935 & 130.5\\
Three Conway summands, forced scan & $>20$ & 0.012196 & $>1639$\\
Ordering only, chain of 64 & 0.014000 & 0.001749 & 8.0\\
Ordering only, chain of 256 & 0.217324 & 0.007209 & 30.1\\
Ordering only, chain of 1024 & 3.414980 & 0.032857 & 103.9\\
\bottomrule
\end{tabular}
\caption{The forced-scan recognition rows disable Alexander in both versions
and Jones in the new version, retaining ordinary reductions and the new
factorization. Ordering rows time only crossing-order selection.}
\label{tab:structural}
\end{table}

The two connected-sum rows demonstrate a structural shortcut independently
of the Jones filter. The old code calculates a product-sized homology,
whereas the new code rejects after the first summand. The twelve-summand
factored-rank run takes 0.031764 seconds and computes two distinct normalized
PD leaves before reusing them; there are twelve leaves altogether. It returns
the full degree convolution as well as $33^{12}$.

The ordering examples are the closures of
$\sigma_1\sigma_2\cdots\sigma_n$ on $n+1$ strands. Both routines return the
same chosen order. The increasing ratios are consistent with the proven
change from quadratic to near-linear ordering work. They say nothing by
themselves about full recognition time on a 1024-crossing arbitrary knot.

\section{Validation and limitations of the evidence}
\label{sec:validation}

The delivered suite contains the original 19 tests and 17 new tests, all
passing. The inherited method-name tests explicitly disable the two new
filters so that they continue to exercise the old pipeline path. This is
documented in the test source; no failing old assertion was silently deleted.
The final recorded run with \code{PYTHONHASHSEED=0} reports 36 tests passing
in 13.120 seconds.

The new tests are mostly property loops rather than single examples. They
compare the heap ordering with the original ordering for every start on 100
diagrams; compare the linear descending predicate with the original one on
200; compare both pivot policies with an independent dense reduced
cube-of-resolutions implementation on 100; and compare the default recognizer
with that dense reference on another deterministic set of 120 generated
diagrams. Small braid closures are restricted to one component by a separate
permutation test.

For the scalar filter, 160 small diagrams are checked against an independent
$2^n$ state sum at two values of $A$, with varied crossing orders. Other tests
check that known unknots are never rejected, that legal row rotations preserve
writhe and invariants, and that bad ring specializations, malformed scan
orders, non-finite time budgets, and undersized work caps are handled
correctly. A composite admissible modulus is tested, so correctness does
not tacitly depend on primality.

Algebra tests compare memoized composition with the retained original
routine and test mutable-alias protection. All 128 three-arc units are
verified to square to the identity. Connected-sum tests compare raw scanning,
factored degree convolution, and the independent dense reference on nine
pairs of small diagrams; cut replay is checked and a tampered cut is rejected.
A twelve-Conway rank test checks exact large-integer output. The difficult
36-crossing scan is run with $d^2=0$ checking enabled.

In addition to the suite, the independent dense reference computes reduced
rank 33 for each of Conway and Kinoshita--Terasaka, in approximately 0.23 and
0.20 seconds. This cross-check is feasible despite their awkwardness for the
old filter pipeline: eleven crossings are small enough for a reference cube.
Its raw outputs, Alexander values, and reducing-move counts are preserved.

\paragraph{What this does not establish.}
There is no Lean formalization of these changes, no proof of correctness of
the Python runtime, and no exhaustive validation over all diagrams. The
36-crossing rank was checked by original versus optimized local algebra and
by chain-complex consistency, not by an independent dense enumeration of
its $2^{36}$ states. Property tests share the inherited PD parser and, in some
cases, its traversal routines. Numerical agreement therefore does not remove
all correlated-error risks. These limits are stated so that the test count
is not mistaken for a proof certificate.

\section{Interfaces, evidence, resource budgets, and reproducibility}
\label{sec:usage}

\subsection{Commands and library entry points}

The package runs directly from the extracted deliverable:
\begin{lstlisting}
python -m fastunknot recognize examples/conway.json
python -m fastunknot recognize examples/hard_unknot_8.json
python -m fastunknot khovanov examples/random5_36.json --check-d2
python -m fastunknot khovanov examples/conway_sum_12.json --factor
python -m unittest discover -s tests -v
\end{lstlisting}
The third command intentionally bypasses simplification and invariant tests.
The fourth explicitly opts into factored ranks. Keeping both interfaces
prevents confusion between a faster full-rank scanner and a shortcut for
recognition or rank multiplication.

The public constructors accept PD, braid, and the inherited rectangular-grid
formats. For example:
\begin{lstlisting}[language=Python]
from fastunknot import Diagram, recognize
D = Diagram.from_braid(3, [1, -2, 1, -2])
answer = recognize(D, seconds=30, max_objects=200_000)
print(answer.status)
print(answer.to_json())
\end{lstlisting}
The raw \code{khovanov_rank(pd)} function assumes an already validated PD;
normal use goes through \code{Diagram.from_pd} or another checked constructor.
The dataclass itself must not be used to bypass those checks.

The old recognition options remain available. New ones are
\code{use_jones}, \code{use_factorization}, the two Jones work caps, and
\code{pivot_strategy}. The corresponding CLI flags include
\code{--no-jones}, \code{--no-factor}, and \code{--pivot lifo}. A completed
mathematical decision has exit code zero, a resource-limited unknown has
code three, and invalid input or a rejected witness has code two.

\subsection{Evidence is contextual, not a universal certificate}

A standalone Jones result can be reproduced on the same diagram:
\begin{lstlisting}
python -m fastunknot jones examples/conway.json > witness.json
python -m fastunknot verify-jones examples/conway.json witness.json
\end{lstlisting}
The verifier recomputes the finite-ring evaluation and compares the relevant
fields; it does not trust the claimed residue. This is a reproducibility
checker, not a succinct polynomial-time checker for a short certificate.
On adversarial inputs it needs external process limits.

The cut-tree replay function reconstructs the leaf diagrams from explicit
two-edge-cut records. A full pipeline result has additional context:
initial R1/R2 traces precede the cuts, and each factor can have its own
reduction trace. Thus a leaf's Jones witness refers to its \emph{post-reduction
leaf}, not necessarily to the original input PD. Those reductions must be
replayed before the respective witness is checked. The current package does
not offer a complete independently implemented verifier for every pipeline
result. Scanner rank output is computational evidence, not a spanning disk
or a formally checked homology certificate.

\subsection{Limits are explicit but cooperative}

\code{max_objects} constrains only the scanner's planned delooped object
count. An exact Jones or Alexander rejection is still valid even if that
ceiling is tiny. A Jones-specific state or transition cap skips Jones and
falls back; exhaustion of the shared time budget stops with \UNKNOWN.
Invalid negative or non-finite budgets are rejected instead of being used as
hidden unlimited modes.

Deadline checks are inserted in reduction, descent, factorization, polynomial
elimination, state transitions, object construction, and cancellation loops.
They are not a hard interruption mechanism: a single large Python arithmetic
or composition operation, or a debug $d^2$ check, may run past its deadline.
Similarly, caches are bounded in entry count but not total bytes, and the
scanner ceiling does not bound morphism support sizes or Python overhead.
Applications requiring strict wall-clock or memory isolation should run the
process with operating-system limits. \code{MemoryError} in the high-level
recognizer becomes \UNKNOWN, but graceful recovery cannot be guaranteed after
an external out-of-memory kill.

\subsection{Rebuilding the measurements and article}

The unchanged comparison code, worker, example generator, and raw results are
included. Run
\begin{lstlisting}
python tools/benchmark.py --repeats 3 --cap 20
python tools/summarize_results.py
\end{lstlisting}
to regenerate the measurements and readable summary. The first command takes
several minutes because it intentionally waits on capped baseline cases.
\code{benchmark_worker.py} exposes the individual operation and ablation
variants for targeted reruns. Fresh processes include newly empty caches in
both variants, and input construction remains outside both timers.

The article has no external figure or bibliography dependencies. From its
\code{docs/} directory, two runs of
\begin{lstlisting}
pdflatex -interaction=nonstopmode -halt-on-error report.tex
\end{lstlisting}
rebuild its PDF and cross-references. LaTeX is not needed to use the software.
The original license is retained, and third-party papers and slides are
cited rather than redistributed in this ZIP.

\section{Conclusion and the next concrete bottleneck}

The delivered improvement is a working exact recognizer, not an unimplemented
proposal. It changes the scanner's elimination behavior enough to complete a
formerly capped raw scan, adds an independent bounded rejection test that
helps the two eleven-crossing Alexander-trivial fixtures, and avoids an
exponential tensor-product basis on exposed connected sums. The strongest
proved bound is \eqref{eq:parameter}, not an unrestricted quasi-polynomial
bound.

For difficult diagrams with no exposed factors and with filters that return
one, the remaining bottleneck is a genuinely large categorical chain complex.
A next implementation step would be to represent repeated boundary matching
multiplicities by blocks and seek algebraically certified low-rank
factorizations of the differential blocks. Such a representation must
preserve the Schur updates and account for the cost of operations on those
blocks; merely storing multiplicities as integers is insufficient when
nontrivial differentials couple the copies. For the separate hierarchy
approach, a compressed topological state and costed implementations of its
moves are still required. Neither direction is smuggled into the present
complexity claim.

\clearpage
\begin{thebibliography}{9}

\bibitem{BN}
D. Bar-Natan,
\emph{Fast Khovanov homology computations},
Journal of Knot Theory and Its Ramifications \textbf{16} (2007), 243--255.
\href{https://arxiv.org/abs/math/0606318}{arXiv:math/0606318}.
Author's publication page:
\href{https://www.math.toronto.edu/drorbn/papers/FastKh/}{FastKh}.

\bibitem{KM}
P. B. Kronheimer and T. S. Mrowka,
\emph{Khovanov homology is an unknot-detector},
Publications Math\'ematiques de l'IH\'ES \textbf{113} (2011), 97--208.
\href{https://arxiv.org/abs/1005.4346}{arXiv:1005.4346}.

\bibitem{Kauffman}
L. H. Kauffman,
\emph{State models and the Jones polynomial},
Topology \textbf{26} (1987), 395--407.
\href{https://doi.org/10.1016/0040-9383(87)90009-7}{doi:10.1016/0040-9383(87)90009-7}.

\bibitem{Ltalk}
M. Lackenby,
\emph{Unknot recognition in quasi-polynomial time},
February 2021, 109-page talk slides.
The exact version supplied in \code{docs/quasipolynomial-talk.pdf}.
\href{https://people.maths.ox.ac.uk/lackenby/quasipolynomial-talk.pdf}{Author-hosted February slides}.
Physical pages 13, 60--74, and 109 are relevant here.

\bibitem{Loxford}
M. Lackenby,
\emph{Unknot recognition in quasi-polynomial time},
March 2021 Oxford slide variant, 109 pages.
\href{https://people.maths.ox.ac.uk/lackenby/quasipolynomial-talk-oxford.pdf}{Author-hosted Oxford variant}.
Cited only to distinguish its stated exponent from the submitted February file.

\bibitem{L2026}
M. Lackenby,
\emph{Incompressible surfaces, hierarchies and unknot recognition},
25 July 2026, version 1, 54 pages.
\href{https://arxiv.org/abs/2607.23350v1}{arXiv:2607.23350v1}.
In particular Section 9, physical pages 34--35.
Public version checked on 18 September 2026.

\bibitem{baseline}
Knots contributors,
\emph{fastunknot: exact unknot recognition with a scanning Khovanov backend},
user-supplied \code{Knots.zip}, \code{fast/} directory, September 2026.
The actual comparison code is retained in \code{baseline/fastunknot/}; source
provenance and the original license accompany this report.

\end{thebibliography}
\end{document}
